\documentclass{situationist-preprint}

\title{The Beginning of Time}
\msid{TSM-14-a}
\msfamily{TSM-14}
\msdiscipline{Time, Space and Modality}
\mstype{Principal manuscript}
\mskeywords{beginning of time; cosmology; singularity theorems; finite past; Kant's first antinomy; Aristotle; cosmological argument}

\begin{document}
\maketitle

\begin{abstract}
Did time have a beginning, and if so, what came before it? We argue that time had a beginning and that the question of what preceded it is ill-formed. Time is the temporal structure of spacetime, and to have a beginning is for that structure to have a finite past: a bound on the duration of every past-directed history, whether or not there is a first moment. The singularity theorems of general relativity, and the theorem of Borde, Guth and Vilenkin for expanding spacetimes, give strong evidence that our universe has a finite past, and the models that evade them either face the problem of accumulating entropy or have a low-entropy boundary from which directed time runs away on both sides. Given that precedence is the earlier-than relation of the temporal series, a question about what preceded every time presupposes a time outside the series and has no answer, not an unknown one. The classical arguments of Aristotle and Kant to the contrary assume a time independent of the world. Whether the universe depends non-temporally on something is a separate question, which the result leaves open.
\end{abstract}

\section{The problem}

Augustine was asked what God was doing before he made heaven and earth, and answered that there was no \enquote{then} before time, since time was itself made \citep[XI.13]{AugustineConfessions}. \citet[VIII.1]{AristotlePhysics} argued the opposite: every moment is a \enquote{now} that is the end of a past and the beginning of a future, so that there can be no first moment, and time and motion are eternal. \citet[A426/B454]{KantCPR} gave proofs both that the world had a beginning in time and that it did not, and took the antinomy to show that the question exceeds possible experience. In the twentieth century the question passed in part to physics. The expansion of the universe, the cosmic background radiation and the singularity theorems of \citet{HawkingPenrose1970} made a finite past the default of relativistic cosmology, and \citet{BordeGuthVilenkin2003} extended the result to inflationary spacetimes. Philosophers have disputed both what the physics shows \citep{Earman1995} and what it would mean: \citet{Craig1979} takes a beginning to require a cause, and \citet{Grunbaum1989} argues that the question of what caused the beginning is a pseudo-problem.

We argue that time had a beginning, and that the question of what preceded it is ill-formed.

\section{Framework}

\begin{definition}[Beginning]\label{d:begin}
Time has a \emph{beginning} if there is a finite bound on the proper time along every past-directed timelike curve from any event, or, in a model with a cosmic time, if the cosmic time is bounded below. A beginning is \emph{closed} if there is a first moment and \emph{open} if there is none, every moment having earlier ones although the past is finite.
\end{definition}

Definition~\ref{d:begin} is stated in terms of proper time along curves, so it does not depend on a global present, which TSM-16 denies. It counts as a beginning the initial singularity of the standard cosmological models, which is not itself a moment of spacetime, so that the beginning is open.

\begin{principle}[Time is temporal structure]\label{p:struct}
Time is the temporal structure of spacetime: the B-series of events ordered by the relation of being earlier than (TSM-01), which is the causal and metrical structure of spacetime. There are no times outside this structure.
\end{principle}

\begin{principle}[Precedence]\label{p:prec}
For $x$ to \emph{precede} $y$ is for $x$ to be earlier than $y$ in the temporal series. Precedence is defined only between elements of that series.
\end{principle}

Principle~\ref{p:struct} is consistent with TSM-17, which argues that there can be time without change: Shoemaker's frozen intervals are stretches of temporal structure in which nothing changes, not times outside the structure. It is also consistent with TSM-06, on which spacetime is a substance; a beginning of time is then a past boundary of that substance, not an event within a larger time.

\section{The argument}

\begin{lemma}[Evidence for a finite past]\label{l:phys}
Our best physical evidence supports a finite past in the sense of Definition~\ref{d:begin}.
\end{lemma}

\begin{proof}
The universe is expanding, and in general relativity, under the energy conditions satisfied by ordinary matter and radiation and under conditions that the observed universe meets, the singularity theorems show that past-directed timelike geodesics are incomplete \citep{HawkingPenrose1970}. The theorem of \citet{BordeGuthVilenkin2003} requires no energy condition: any spacetime whose average expansion rate along a past-directed geodesic is positive is past-incomplete, which covers eternally inflating models. The remaining candidates for an infinite past are models in which a contracting phase precedes ours. Cyclic models of this kind, as \citet{Tolman1934} observed, accumulate entropy from cycle to cycle, so that an infinite past of cycles is incompatible with our low-entropy past; and in models with a single bounce or a minimum of entropy, such as those of \citet{CarrollChen2004} and \citet{BarbourKoslowskiMercati2014}, entropy increases away from the minimum on both sides. By TSM-05, time has a direction because of a low-entropy boundary condition, and the direction points away from the low-entropy boundary, here on each side of the minimum. The other side is not earlier than ours in the directed sense; its own past begins at the same boundary. In each of these models, then, directed time has a beginning.
\end{proof}

\begin{lemma}[Ill-formedness]\label{l:ill}
The question \enquote{What preceded the beginning of time?} has no answer, not even a negative one of the form \enquote{nothing preceded it, but something could have}.
\end{lemma}

\begin{proof}
By Principle~\ref{p:prec}, an answer would be something earlier than the earliest times, or, if the beginning is open, earlier than all times. By Principle~\ref{p:prec}, being earlier is a relation between elements of the series, and by Principle~\ref{p:struct} there are no times outside it, so whatever is earlier than a time is itself in the series and is not earlier than all its elements. The question presupposes that there is a position in the temporal order before all positions, and the presupposition is contradictory. A question with a contradictory presupposition is ill-formed. Nor is \enquote{nothing} a correct answer in the sense of reporting an empty period: an empty period would be a stretch of temporal structure, and therefore part of time.
\end{proof}

\begin{theorem}\label{t:main}
Time had a beginning, and the question of what preceded it is ill-formed.
\end{theorem}

\begin{proof}
Lemmas~\ref{l:phys} and~\ref{l:ill}.
\end{proof}

The first conjunct is empirical and rests on the evidence of Lemma~\ref{l:phys}; the second is conceptual and holds of any finite past. \citet{Hawking1988} compared the question to asking what lies north of the North Pole: in the no-boundary proposal the past is finite without a first moment, and asking what came earlier is like asking for a point further north.

\section{Objections and replies}

\begin{objection}[Aristotle]\label{o:arist}
Every moment is the end of a past. A first moment would be the end of no past, which is impossible; so time had no beginning \citep[VIII.1]{AristotlePhysics}.
\end{objection}

\begin{reply}
The objection, if sound, excludes a closed beginning, not an open one. On an open beginning every moment has earlier moments, as Aristotle requires, but the earlier moments are confined to a finite interval, as the real numbers greater than zero have no least element and are bounded below. The beginning that the singularity theorems indicate is of this kind.
\end{reply}

\begin{objection}[Kant]\label{o:kant}
If the world had a beginning, there was an empty time before it, and in empty time no moment is distinguished from another, so nothing could have determined the world to begin when it did \citep[A427/B455]{KantCPR}.
\end{objection}

\begin{reply}
The argument assumes a time independent of the world, within which the world begins. Principle~\ref{p:struct} denies the assumption: time is the temporal structure of spacetime, and by Lemma~\ref{l:ill} an empty time before the beginning would be part of time and so not before it. Leibniz made the same reply to Clarke's question why God did not create the world sooner: there was no sooner \citep{LeibnizClarke1716}.
\end{reply}

\begin{objection}[What begins has a cause]\label{o:kalam}
Whatever begins to exist has a cause of its existence. If the universe began, it has a cause, and the question of what brought it about is not ill-formed \citep{Craig1979}.
\end{objection}

\begin{reply}
Theorem~\ref{t:main} concerns precedence, not dependence. That nothing preceded the beginning does not settle whether the universe, with its time, depends non-temporally on something, and the question whether it does is well-formed. It is the question of MET-01-a, which answers why there is something by the structure of the space of possibilities rather than by a prior event, and of REL-01, which argues that cosmological considerations converge on a necessary being but leave its causal role undetermined. What the objection cannot have is a cause that acts before the beginning; a cause of the universe, if there is one, stands to it as a ground or a timeless agent does, not as an earlier event does. \citet{Grunbaum1989} is right that the temporal version of the question of creation is a pseudo-problem; he is not thereby right about the non-temporal one, on which we take no stand.
\end{reply}

\begin{objection}[Quantum gravity]\label{o:qg}
The singularity theorems belong to classical general relativity, which fails near the singularity. A quantum theory of gravity may replace the singularity with a bounce and an infinite past.
\end{objection}

\begin{reply}
The objection is the strongest, and it shows that the first conjunct of Theorem~\ref{t:main} is empirical. But it is not evidence against a beginning. The Borde--Guth--Vilenkin result does not use the classical field equations, only the geometry of expansion, so a bounce must reverse the expansion; and in bounce models with a low-entropy bounce the argument of Lemma~\ref{l:phys} applies, so that directed time begins at the bounce. If quantum gravity yields a model with an infinite directed past, the first conjunct will have been refuted; no such model now has evidential support comparable to the finite-past models.
\end{reply}

\section{Conclusion}

Time is the temporal structure of spacetime, and the evidence of relativistic cosmology supports a finite past, open rather than closed: every moment has earlier ones, but none lies more than a finite duration back. Since precedence is a relation within the series, the question of what preceded every time presupposes a time outside the series and has no answer. Aristotle's argument excludes only a first moment, and Kant's assumes a time independent of the world. Whether the universe depends on something non-temporally remains a well-formed question, and nothing here answers it.

\begin{thebibliography}{99}\small
\bibitem[Aristotle(Phys.)]{AristotlePhysics} Aristotle. \emph{Physics}. Trans. R.~P. Hardie and R.~K. Gaye, in J.~Barnes (ed.), \emph{The Complete Works of Aristotle}, Princeton University Press, 1984.
\bibitem[Augustine(Conf.)]{AugustineConfessions} Augustine. \emph{Confessions}. Trans. H.~Chadwick, Oxford University Press, 1991.
\bibitem[Barbour et~al.(2014)]{BarbourKoslowskiMercati2014} Barbour, J., Koslowski, T., and Mercati, F. (2014). Identification of a gravitational arrow of time. \emph{Physical Review Letters}, 113, 181101.
\bibitem[Borde et~al.(2003)]{BordeGuthVilenkin2003} Borde, A., Guth, A.~H., and Vilenkin, A. (2003). Inflationary spacetimes are incomplete in past directions. \emph{Physical Review Letters}, 90, 151301.
\bibitem[Carroll and Chen(2004)]{CarrollChen2004} Carroll, S.~M., and Chen, J. (2004). Spontaneous inflation and the origin of the arrow of time. arXiv:hep-th/0410270.
\bibitem[Craig(1979)]{Craig1979} Craig, W.~L. (1979). \emph{The Kal\=am Cosmological Argument}. Macmillan.
\bibitem[Earman(1995)]{Earman1995} Earman, J. (1995). \emph{Bangs, Crunches, Whimpers, and Shrieks: Singularities and Acausalities in Relativistic Spacetimes}. Oxford University Press.
\bibitem[Gr\"unbaum(1989)]{Grunbaum1989} Gr\"unbaum, A. (1989). The pseudo-problem of creation in physical cosmology. \emph{Philosophy of Science}, 56(3), 373--394.
\bibitem[Hawking(1988)]{Hawking1988} Hawking, S.~W. (1988). \emph{A Brief History of Time}. Bantam.
\bibitem[Hawking and Penrose(1970)]{HawkingPenrose1970} Hawking, S.~W., and Penrose, R. (1970). The singularities of gravitational collapse and cosmology. \emph{Proceedings of the Royal Society of London A}, 314, 529--548.
\bibitem[Kant(1781)]{KantCPR} Kant, I. (1781/1787). \emph{Critique of Pure Reason}. Trans. P.~Guyer and A.~W. Wood, Cambridge University Press, 1998.
\bibitem[Leibniz and Clarke(1716)]{LeibnizClarke1716} Leibniz, G.~W., and Clarke, S. (1715--1716). \emph{The Leibniz--Clarke Correspondence}. Ed. H.~G. Alexander, Manchester University Press, 1956.
\bibitem[Tolman(1934)]{Tolman1934} Tolman, R.~C. (1934). \emph{Relativity, Thermodynamics and Cosmology}. Clarendon Press.
\end{thebibliography}

\end{document}
